\section{Introduction}
\label{sec:introduction}

Code agents use large language models to solve software engineering tasks through repeated interaction with external tools, such as the file system and shell \citep{yao2023react,yang2024sweagent}. This closed-loop process allows an agent to inspect a repository, modify code, and validate its solution through execution feedback. However, it also makes agent inference expensive.

One important reason is that large language models are typically post-trained to prioritize task success rather than execution efficiency. Their rewards encourage correct solutions but rarely penalize redundant exploration, unnecessary tool calls, or excessively long tool outputs. As a result, a code agent may repeatedly inspect similar files or execute uninformative commands as long as it eventually solves the task. On repository-level benchmarks such as Deep-SWE, a single issue may therefore require hundreds of interaction rounds. Each round adds actions and observations to the trajectory, namely the sequence of agent actions and environment feedback generated during execution. These observations remain in the context of later model calls, so inefficient behavior increases both the number of calls and the cost of each call.

Existing approaches mainly reduce this cost by compressing the accumulated interaction history. When the context becomes large, they summarize earlier messages or remove information considered unimportant. Although this shortens later prompts, it acts only after costly interactions have already occurred. It also introduces two practical risks. First, rewriting earlier messages breaks the stable prompt prefix used by prefix caching, which reuses previously processed context across consecutive calls. Second, observations that appear unnecessary at one stage may become important later for locating the root cause of a bug. History compression must therefore balance token reduction against information loss, while leaving the underlying inefficient behavior unchanged. A more direct solution is to prevent unnecessary interactions from being generated in the first place.

The behavior of a code agent is shaped not only by its model parameters but also by its execution harness. The harness is the surrounding system that defines the available tools, tool interfaces, execution rules, observation formats, and instructions used during agent execution. Modifying these components can guide the behavior of an existing model without updating its parameters. A cost-aware harness could reduce redundant inspections, limit uninformative outputs, and encourage more efficient debugging strategies. Harness optimization therefore offers a natural way to adapt an accuracy-oriented model to cost-sensitive deployment.

However, manually designing such a harness is difficult. Expensive behavior is often model-specific and hidden within long execution trajectories. The same action may be necessary in one context but redundant in another, making general cost patterns difficult to identify without analyzing many historical executions. Even after an expensive pattern is discovered, the appropriate harness change is not obvious. Restricting file inspection may reduce observation length but increase the number of tool calls, while truncating outputs may remove information required for debugging. Because tools, execution rules, and instructions jointly affect agent behavior, an apparently reasonable modification may introduce new inefficiencies or reduce task performance.

These observations suggest using prior executions to change the mechanism that produces future interactions, rather than rewriting the history of an active rollout. We develop \sys{}, a trajectory-to-harness optimization framework for this purpose. The framework keeps the model fixed and treats the tool registry, execution dispatcher, and behavioral instruction as a coupled update object. During each rollout, this harness remains fixed and the interaction history grows only by appending new messages. Adaptation occurs between rollouts.

The main difficulty is turning a long execution trace into evidence for a reusable harness change. \sys{} first annotates each action using only the trajectory prefix available at that action, producing a predicted dependency graph together with operation-type and execution-state labels. It then extracts bounded records around several cost-relevant patterns, such as actions outside the predicted final dependency closure, mergeable sibling operations, repeated failures before a successful pivot, and unnecessarily long or repeated observations. Each record pairs an observed local slice with a proposed compact view. The proposed view is a diagnostic hypothesis rather than an executed improvement.

These records guide a unified update of the three harness artifacts. Because unrestricted self-editing can add permanent schema overhead or introduce inconsistent tool behavior, every candidate is edited in staged state and passes two static screens before internal acceptance. Deterministic checks enforce the artifact contracts, and a separate language-model judge assesses generality, prospective cost, correctness, consistency, and alignment with the supplied evidence. A rejected batch restores the preceding staged harness. This gate is deliberately fail-closed, but it does not execute the candidate or establish that task utility is preserved. A completed folded harness is tested only when it drives a later rollout. Figure~\ref{fig:overview} summarizes this loop.

By changing the harness used for later executions, \sys{} targets the source of avoidable interactions while leaving each accumulated rollout history intact. This makes the intervention complementary to context compression and compatible with exact-prefix caching when the provider recognizes the unchanged prefix. It does not assume that cached tokens are free, nor does it claim that the predicted dependencies are causal.

Our contributions are summarized as follows:

\begin{itemize}[leftmargin=*,itemsep=2pt,topsep=3pt]

\item We formulate cost-aware trajectory-to-harness optimization as an upstream alternative to rewriting an active interaction history. The formulation makes the amortized adaptation cost and downstream utility constraint explicit.

\item We develop the first framework that turns dependency-annotated code-agent trajectories into gated, unified harness updates specifically for inference-cost reduction while leaving each rollout's accumulated history append-only.

\item We introduce a concrete update process based on future-suffix-blind action annotation, bounded graph and behavioral records, coupled edits to tool schemas, dispatch semantics, and instructions, and fail-closed static screening. We separate this provisional screening from execution-based evidence of cost or utility.

\end{itemize}
